\chap[Appendices]{Appendix A. Key Formulas}
The formulas a technician uses, in plain text for quick reference. The chapter gives the full statement and a worked example.

\sect{Signals, points and links (Chapters 2 to 5)}
\begin{tbl}{P{0.4\textwidth}L}
Law & Formula\\\midrule
2.1 Sense resistor & V = I × R; 4–20 mA × 250 Ω = 1–5 V\\
2.2 Loop budget & V\_supply ≥ V\_tx,min + I\_max × (R\_sense + R\_cable + R\_other)\\
2.3 Scaling & EU = EU\_min + (raw − raw\_min) / (raw\_max − raw\_min) × (EU\_max − EU\_min)\\
2.4 Resolution & Δ = span / 2\^{}n\\
2.5 Nyquist & f\_s > 2 × f\_max (in practice 5 to 10 ×)\\
2.6 Counter delta & ΔC = (C\_new − C\_old) mod 2\^{}n\\
2.7 Pulse flow & Q = K × ΔC / Δt\\
3.1 Deadband & report if |x − x\_last| > d\\
3.2 Ordering events & Δ > e\_1 + e\_2\\
3.3 Drift & e = e\_0 + y × t; 1 ppm = 86.4 ms a day\\
4.1 Serial byte time & t = 10 / baud\\
4.2 Poll cycle & T = Σ (request + reply + turnaround)\\
4.3 Delivery with retries & P = 1 − p\^{}k\\
4.4 Data volume & R = 8 × λ × B; D\_month = λ × B × 86,400 × 30\\
4.5 Latency & T\_total = sum of all stage delays\\
5.1 Fan-out & λ\_out = s × λ\\
5.4 Lost Sparkplug messages & m = (received − expected) mod 256\\
5.5 Outage buffer & B = λ × b × T\_out\\
6.1 OPC UA queue size & q ≥ ceiling(T\_p / T\_s)\\
6.3 Edge aggregation & ratio = T\_a / T\_s\\
7.1 GOOSE repeats & T\_k = T\_min × 2\^{}k until T\_k ≥ T\_max\\
8.1 Clock offset & θ = ((t\_2 − t\_1) + (t\_3 − t\_4)) / 2; δ = (t\_4 − t\_1) − (t\_3 − t\_2)\\
8.2 Asymmetry error & ε = (d\_out − d\_back) / 2\\
8.3 Recovery loss & N\_lost = f × T\_rec\\
\end{tbl}

\sect{The SCADA system (Chapters 9 to 14)}
\begin{tbl}{P{0.4\textwidth}L}
Law & Formula\\\midrule
9.1 Update rate & λ = c × N\\
9.2 History storage & S = (N\_h / T\_log) × b × T\_ret\\
9.3 Amdahl & S(n) = 1 / ((1 − p) + p / n)\\
9.4 Queue delay & ρ = λ / μ; W = 1 / (μ − λ)\\
9.5 Availability & A = MTBF / (MTBF + MTTR)\\
9.6 Parallel and series & A\_par = 1 − (1 − A)\^{}n; A\_ser = A\_1 × A\_2 × …\\
10.2 Alarm rate & r = N\_alarms / (T × N\_operators)\\
10.3 Time in flood & F = flood intervals / total intervals\\
10.4 Bad-actor share & S\_k = top k counts / all counts\\
10.5 Alarm deadband & clear when x < L − d; d ≈ 6σ or more\\
10.6 Asset tree & N\_leaf = b\^{}L\\
11.3 Interpolation & v(t) = v\_1 + (v\_2 − v\_1) × (t − t\_1) / (t\_2 − t\_1)\\
11.4 Time-weighted average & v\_avg = Σ(v\_i × Δt\_i) / ΣΔt\_i\\
11.5 Totalising & V ≈ Σ (Q\_i + Q\_i+1) / 2 × Δt\\
11.6 Interface count & I\_p2p = n(n − 1) / 2; I\_UNS = n\\
12.2 First-order response & y(t) = K × Δu × (1 − e\^{}(−(t − θ)/τ))\\
12.3 Lambda tuning & K\_c = τ / (K × (λ + θ)); T\_i = τ\\
12.6 Limit shift & Δ = z × (σ\_1 − σ\_2)\\
12.7 Affinity laws & Q ∝ N; H ∝ N²; P ∝ N³\\
13.1 Break-even & n* = t\_dev / (t\_m − t\_a)\\
13.2 Error accumulation & P = 1 − (1 − p)\^{}n\\
13.4 Change downtime & T\_down = n × CFR × MTTR\\
14.1 Scan response & T\_worst = T\_in + 2 T\_scan + T\_out\\
14.3 Precision & s π / (s π + f (1 − π))\\
\end{tbl}

\sect{Security and practice (Chapters 15 to 20)}
\begin{tbl}{P{0.4\textwidth}L}
Law & Formula\\\midrule
15.1 Brute force & T = 2\^{}n / r\\
15.2 Entropy & H = L × log2(S)\\
15.4 Renewals & N / V per year\\
16.1 Trust entries & E\_self = 2cs; E\_CA = c + s\\
16.2 Key count & K\_pairs = n(n − 1) / 2; K\_star = n\\
17.1 Possible conduits & C\_max = z(z − 1) / 2\\
17.2 Security gap & g\_i = max(0, SL-T\_i − SL-A\_i)\\
17.3 Dependency & A\_action = A\_local × A\_remote\\
18.1 Detection & P = 1 − (1 − d)\^{}T\\
18.2 Barriers & P = p\_1 × p\_2 × …\\
19.2 Exposure & P = 1 − e\^{}(−hW)\\
19.3 SOCI reporting & 12 h significant; 72 h relevant; 84 h written record after verbal\\
20.1 Half-splitting & k = ceiling(log2 N)\\
\end{tbl}

\sect{Numbers worth memorising}
\begin{tbl}{P{0.4\textwidth}L}
Quantity & Value\\\midrule
Live zero; fault limits & 4 mA; below about 3.6 or above about 21 mA\\
Seconds in a day; a year & 86,400; 31,536,000\\
2\^{}8, 2\^{}16, 2\^{}32 & 256; 65,536; 4,294,967,296\\
Copper resistivity & 0.0175 Ω·mm²/m\\
1 − 1/e & 63.2\%\\
Ports & Modbus 502; secure Modbus 802; IEC 104 2404; OPC UA 4840; MQTT 1883 plain, 8883 TLS\\
Sparkplug sequence range & 0 to 255\\
Alarm benchmarks & about 1 per 10 minutes comfortable; flood above 10 in 10 minutes\\
z for 95\% one-sided & 1.645\\
\end{tbl}

\chap{Appendix B. Study Plan and Self-Test}
\sect{B.1 A ten-week plan}
About five hours a week: read the chapters, work every example by hand, do the exercises without looking back, then the practical task.
\begin{tbl}{P{0.1\textwidth}P{0.14\textwidth}L}
Week & Chapters & Practical task\\\midrule
1 & 1, 2 & Take a real loop on site. Measure the current, calculate the engineering value by hand, compare with the screen. Draw the chain of Section 1.3 for one point\\
2 & 3, 4 & Find where your timestamps are applied and where your clocks get their time. Capture a Modbus or DNP3 exchange in Wireshark and identify each field\\
3 & 5, 6 & Install a broker and MQTT client; publish, subscribe, test wildcards and Last Will. Browse an OPC UA server with a generic client; find a queue size and the security mode of every endpoint\\
4 & 7, 8 & Read an SCL file. List your network's redundancy method and the longest gap each application tolerates\\
5 & 9, 10 & Find your system's point count and server specification; estimate storage per year. Export a week of alarms and compute Laws 10.2 to 10.4\\
6 & 11, 12 & Pull history for one flow and total it by hand. Step-test a loop (with permission) and calculate K, θ and τ\\
7 & 13, 14 & Write the functions for Laws 2.3, 2.6 and 5.4 in Python. Put a configuration export under Git, change it, and read the diff\\
8 & 15, 16 & Inspect a real certificate: issuer, expiry, key size. Check which protocols on your system run with no security, and why\\
9 & 17, 18, 19 & Draw your zones and conduits from risk. List every remote-access path. Work Worked Example 19.2 for your own site\\
10 & 20 & Write an ITP for one subsystem. Then fault-find something real using the method of Section 20.3, and write it up\\
\end{tbl}

\sect{B.2 Self-test}
Answer aloud, then check. If you can answer all thirty, you know more than most people in the room.
\begin{enumerate}
\item What three things make up a point? Value, timestamp, quality.
\item Why does a 4–20 mA loop start at 4 mA? So a broken wire (0 mA) differs from a true zero.
\item A reading flattens below the top of its range however high the level goes. What should you check? The loop voltage budget.
\item What is the danger of scaling in two places? The value is scaled twice and is wrong.
\item Why use a normally closed contact for a high-high alarm? A broken wire or lost supply also gives the alarm.
\item Where should a timestamp be applied, and why? At the field device, so order and time survive polling and outages.
\item What are the classic Modbus faults? Off-by-one addressing, word order, scaling.
\item Why does DNP3 deliver correct data after an outage? Events are buffered and timestamped at the outstation.
\item In MQTT, who knows the address of a subscriber? Only the broker.
\item What three things does Sparkplug B add to MQTT? Topic structure, payload format, state management.
\item What must a Sparkplug node send before any data? NBIRTH, listing every metric.
\item Why can a historian reading OPC UA lose values while the display looks fine? The monitored item queue is smaller than T\_p / T\_s.
\item What does BadSubscriptionIdInvalid after a server restart mean? The client must create a new subscription.
\item Why does GOOSE run at layer 2? Trip times of a few milliseconds leave no time for the IP stack.
\item Why can PTP beat NTP by three orders of magnitude? Hardware timestamps and PTP-aware switches.
\item Why does a server at 99\% load feel far worse than one at 75\%? Queue delay explodes near full load.
\item Why are two virtual servers on one host not redundant? They share a failure.
\item When is an indication not an alarm? When there is no defined operator action.
\item Why is a plain average of historian values usually wrong? Exception-reported values must be time-weighted.
\item What does a UNS hold, and what does it not? Current state; history stays in a historian.
\item A loop cycles in manual as well as auto. Is it the tuning? No: look at the process, valve or another loop.
\item What are the three numbers of a first-order process model? Gain, dead time, time constant.
\item Name the four GitOps principles. Declarative, versioned and immutable, pulled automatically, continuously reconciled.
\item Advisory or closed-loop AI: why does it matter? They need different risk assessments.
\item An RTU's TLS fails after its clock reset. Why? The certificate appears not yet valid.
\item What is the main weakness of a shared pre-shared key? One leak exposes every device using it.
\item What is the difference between SL-C and SL-A? What a product can deliver against what was achieved once installed.
\item Name two incidents where no vulnerability was exploited. Any two of Maroochy, Unitronics, FrostyGoop, Colonial.
\item How long does an Australian critical infrastructure entity have to report a critical cyber incident? 12 hours from becoming aware.
\item What is the first step in fault-finding? Define the symptom exactly; then find what changed.
\end{enumerate}

\chap{Appendix C. Glossary and Sources}
\sect{C.1 Glossary}
\begin{tbl}{P{0.2\textwidth}L}
Term & Meaning\\\midrule
AESCSF & Australian Energy Sector Cyber Security Framework\\
APC & Advanced process control\\
APL & Advanced Physical Layer: two-wire Ethernet for process instruments\\
bdSeq & Sparkplug birth-death sequence number\\
Broker & MQTT server that routes messages\\
CA & Certificate authority\\
CIRMP & Critical Infrastructure Risk Management Program\\
CRA & EU Cyber Resilience Act\\
CVSS & Common Vulnerability Scoring System\\
DCOM & Distributed COM: Microsoft's remote object system, used by OPC Classic\\
DMZ & Demilitarised zone between control and business networks\\
DNP3 & Distributed Network Protocol, IEEE 1815\\
DTLS & TLS for UDP\\
FAT, SAT & Factory and Site Acceptance Tests\\
FR & Foundational requirement (IEC 62443)\\
GDS & Global Discovery Server (OPC UA certificates)\\
GOOSE & Generic Object Oriented Substation Event (IEC 61850)\\
HART & Digital signal on a 4–20 mA loop\\
HMI & Human–machine interface: the operator's screens\\
HSR, PRP & Seamless redundancy protocols, IEC 62439-3\\
IED & Intelligent electronic device\\
ITP & Inspection and Test Plan\\
JSON & JavaScript Object Notation\\
MAC & Message authentication code\\
Mimic & A SCADA graphic (Geo SCADA's word)\\
MMS & Manufacturing Message Specification\\
MPC & Model predictive control\\
MQTT & Lightweight publish-subscribe messaging protocol\\
MTP & Module Type Package\\
NTP, PTP & Network Time Protocol; Precision Time Protocol, IEEE 1588\\
OPC DA, UA & OPC Data Access (Classic); OPC Unified Architecture, IEC 62541\\
PLC & Programmable logic controller\\
PSK & Pre-shared key\\
QoS & MQTT quality of service level\\
RBE & Report by exception\\
RTU & Remote terminal unit\\
SBOM & Software bill of materials\\
SCL & Substation Configuration Language\\
SIS & Safety instrumented system\\
SL-T, SL-C, SL-A & Security level target, capability, achieved\\
SOCI & Security of Critical Infrastructure Act 2018\\
SOE & Sequence of events\\
TLS & Transport Layer Security\\
TSN & Time-Sensitive Networking, IEEE 802.1\\
UNS & Unified Namespace\\
VLAN & Virtual LAN\\
\end{tbl}

\sect{C.2 Sources}
Mathematics and protocol rules are from the standards named in each chapter. Product and legal facts were checked against these sources on 16 and 19 November 2026.
\begin{itemize}
\item Geo SCADA Expert releases, 2025 and 2025 R2 notes, and Python articles, Schneider Electric Knowledge Base; Python Geo SCADA Client Samples, GitHub
\item Configure the Sparkplug Publisher; AVEVA PI Server and CONNECT connection parameters, Geo SCADA Expert 2025 help
\item VBScript deprecation: timelines and next steps, Microsoft
\item Cyber Resilience Act reporting obligations, Matheson; EU Cyber Resilience Act: 24-hour reporting duties, Jones Day
\item Report a cyber security incident, ASD's Australian Cyber Security Centre; mandatory cyber incident reporting guidance, and SOCI Act 2018 for water and sewerage, Cyber and Infrastructure Security Centre
\item Cyber Security (Ransomware Payment Reporting) Rules 2025, explanatory statement, Federal Register of Legislation; ransomware payment reporting in force, MinterEllison
\item Principles of operational technology cyber security, ASD's Australian Cyber Security Centre (October 2024)
\item Dragos 2026 OT Cybersecurity Report and Year in Review (press release and summaries)
\item FrostyGoop: Dark Reading, The Register, TechCrunch; the SCADASEC response as summarised by Trout Software
\item Maroochy Water Breach, MITRE ATT\&CK for ICS (C0020); The Register, 31 October 2001
\item Better alarm handling, UK Health and Safety Executive (Milford Haven, 1994); alarm deadband guidance, exida
\end{itemize}

\sect{C.3 Items to verify before teaching others}
\begin{itemize}
\item[!] Product features marked \emph{verify}: Geo SCADA's JSON configuration export, its viewer for AVEVA OMI, and the 1.5 million point figure quoted at the conference
\item[!] The SOCI asset threshold and CIRMP framework level that apply to your own utility
\item[!] Incident figures in Section 18.1 against each primary report (FrostyGoop's impact is disputed)
\item[!] Byte counts and timings in Worked Examples 4.1 and 16.1 against your own equipment
\item[!] Historic record size in bytes, to firm up the storage estimate in Worked Example 9.1
\end{itemize}

\chap{Appendix D. Questions to Ask a Vendor}
A good question shows more understanding than a recited fact, and protects the site from buying a problem. These work on any product.

\sect{D.1 Questions that work on any stand}
\begin{itemize}
\item What happens when it loses connection: to the cloud, to the broker, to the time source?
\item What does the migration look like on a brownfield site, with the plant running?
\item Which standard is it certified against, and which part? (For IEC 62443: 4-1, 4-2, or both?)
\item Who holds the keys and certificates, and what happens when they expire?
\item Can you show it at scale, with numbers from a reference site?
\item Can I export the configuration in a format I can compare and store in Git?
\item What is the worst-case figure, not the average?
\end{itemize}

\sect{D.2 Topic by topic}
\begin{tbl}{P{0.25\textwidth}L}
Topic & Question\\\midrule
MQTT and UNS & Who is authorised to publish commands through the broker? What happens to data integrity when the primary host is down for an hour?\\
OPC UA & Which companion specifications does the server expose natively? How are certificates renewed across hundreds of devices?\\
Field protocols & What is the key rotation process for devices on serial or low-bandwidth links? With TLS on, how does monitoring still see function codes?\\
Alarms & Is state-based suppression native, or custom logic in the PLC? How are shelved alarms reported at handover?\\
Time & What is the holdover when GPS is lost? Is every switch PTP-aware?\\
Virtual PLCs & What is the redundancy design when the controller is a virtual machine? What is the worst-case jitter under load?\\
AI & Is it advising or writing to the controller? What permissions does it hold, and where is the audit trail?\\
Security & Which requirements does the product meet natively, and which need compensating controls?\\
\end{tbl}

\sect{D.3 Twelve principles to carry}
\begin{enumerate}
\item A value without time and quality is not information. (Chapter 3)
\item Scale once, in one agreed place. (Rule 2.1)
\item MQTT is transport, Sparkplug is state, OPC UA is meaning. (Chapters 5, 6)
\item A connected session does not prove data is flowing. (Section 6.3)
\item The weakest series element sets the availability. (Law 9.6)
\item If the operator has no action, it is not an alarm. (Definition 10.1)
\item The data model comes before the broker. (Rule 11.2)
\item Fix the plant before the tuning. (Rule 12.1)
\item Security level is a vector, and SL-C is not SL-A. (Chapter 17)
\item Engineer out the worst consequence physically. (Section 17.4)
\item The clock starts when you become aware. (Law 19.3)
\item Read the log before guessing. (Rule 20.2)
\end{enumerate}
